\documentclass[10pt,a4paper]{amsart}
\usepackage[margin=19mm]{geometry}
\usepackage{amsmath,amssymb,mathtools}
\usepackage[scr=rsfs,cal=euler]{mathalfa}
\usepackage{xfrac}
\usepackage[unicode,hidelinks,bookmarks=false]{hyperref}
\newcommand{\ie}{i.e.\ }
\newcommand{\cf}{cf.\ }
\newcommand{\resp}{respectively\ }
\newcommand{\Subsection}{\subsection}
\providecommand{\nobreakdash}{\nobreak\hbox{-}}
\setlength{\emergencystretch}{2em}
\multlinegap0pt
\usepackage{bm}
\usepackage{bbm}
\usepackage{dsfont}
\usepackage{xspace}
\usepackage{cancel}
\usepackage{mathtools}
\usepackage{amsthm}
\makeatletter
\@ifundefined{lemma}{\newtheorem{lemma}{Lemma}}{}
\makeatother
\newcommand{\RFA}{\mathbb{R}_{\mathcal{F}(A)}}
\newcommand{\RR}{\mathbb{R}}
\title[Fuzzy Variational Calculus in Linearly Correlated Space: Par]{Fuzzy Variational Calculus in Linearly Correlated Space: Part I}
\author{Gast\~ao S. F. Frederico, Estev\~ao Esmi,, La\'ecio C. Barros}
\date{Source: arXiv:2503.07612v1 (CC BY 4.0)}
\begin{document}
\maketitle

\noindent\emph{Source and licence.} Fuzzy Variational Calculus in Linearly Correlated Space: Part I. Version arXiv:2503.07612v1 (CC BY 4.0), distributed under the Creative Commons Attribution 4.0 International licence, as recorded in the arXiv licence metadata for this exact version and shown on the abstract page https://arxiv.org/abs/2503.07612v1. Full rights notice: licenses/arxiv-2503.07612-CC-BY-4.0.txt.\par\smallskip

\noindent\textit{Editorial scope.} Counted unit: the Lemma whose source label is \texttt{lem:fundamental1} in the article (source0.tex lines 796--813, source label \texttt{lem:fundamental1}), delivered together with the article's own complete original proof of it (source0.tex lines 814--880, one \texttt{proof} environment of 67 lines). No mathematics is written, repaired or re-derived: statement and proof are the article's own text line for line. Required context retained uncounted: none. Every definition, display and label of those lines is retained unchanged. Omitted from this package and not redistributed: 1--795, 881--1265. Nothing inside a retained excerpt is omitted or altered: every retained body is delivered exactly as the article's lines are written, so no omission receipt of any kind accompanies this package. Citations: the retained excerpts carry no citation command at all, so no bibliography entry of the article is transcribed and no reference is redistributed. Reference adaptation: none was needed and none was made; every reference inside the delivered unit resolves inside it, so no unresolved reference or citation is produced. Printed slips kept literal: no printed word, symbol, hypothesis or number of the counted statement or of its proof was altered. Layout and class: the article's own class is \texttt{article} with packages including authblk, inputenc, latexsym, amsthm, amsmath, amssymb, mathtools, color, hyperref, graphics, graphicx, subcaption, lineno; the wrapper is the lane's pinned \texttt{amsart} editorial wrapper at 10pt on A4, which restates the theorem-like environments the retained text needs and adds the article's own macro definitions that the retained text needs (\texttt{\textbackslash RFA}, \texttt{\textbackslash RR}), with extra wrapper packages (bm, bbm, dsfont, xspace, cancel, mathtools). The delivered numbering is the unit's own. Assets: no retained span contains a figure environment, a graphics-inclusion command, a PDF-page inclusion, a table or a TikZ picture; no image file is copied into this package, the package declares no asset dependency and no image file is redistributed with it. Licence: version arXiv:2503.07612v1 of this article is distributed under the Creative Commons Attribution 4.0 International licence, as recorded in the arXiv licence metadata for this exact version and shown on the abstract page https://arxiv.org/abs/2503.07612v1. A full rights notice is delivered at licenses/arxiv-2503.07612-CC-BY-4.0.txt. Characters: every retained excerpt is pure ASCII, so the delivered engine, XeLaTeX with its default Latin Modern text font, reports no missing character.
\begin{lemma}(Lagrange)\label{lem:fundamental1}
Let $A$ be an asymmetric fuzzy number such that $[A]_1 = \{a_m\}$. 
\begin{itemize}
\item[(i)] If $f:[a,b]\longrightarrow\RFA$ is an integrable function such that  
\begin{equation}\label{eq:fundamental1a}
\int_{a}^{b} f(t)\odot\eta(t)\textnormal{d}t \in \RFA^0	
\end{equation}
for any $\eta \in C^{l}\left([a,b],\RFA\right)$ with $\eta^{(j)}(a)=\eta^{(j)}(b)=0$, $j=0,1,\dots,l$, then 
$f(t) \in \RFA^0$ almost everywhere in $[a,b].$

 \item[(ii)] If $f:[a,b]\longrightarrow\RFA$ is an integrable function such that  
\begin{equation}\label{eq:fundamental1b}
\int_{a}^{b} f(t)\odot\eta(t)\textnormal{d}t=0	
\end{equation}
for any $\eta \in C^{l}\left([a,b],\RFA\right)$ with $\eta^{(j)}(a)=\eta^{(j)}(b)=0$, $j=0,1,\dots,l$, then 
$f(t) = 0$ almost everywhere in $[a,b].$
\end{itemize}
\end{lemma}

\begin{proof}
First of all, there exist integrable functions $r,q:[a,b]\to \RR$ such that $f(t) = r(t) + q(t)A$ for all $t \in [a,b]$. Since $r$ and $q$ are integrable functions, they are continuous almost everywhere in $[a,b]$, that is, they are discontinuous in subsets with zero measure. Thus, there exist a open set $I \subset [a,b]$ such that $r$ and $q$ are continuous in $I$ and  the closure of $I$ is $[a,b]$. 
Let $\epsilon> 0$, consider the functions $\delta_k:[-\epsilon,\epsilon]\to \RR$ given by 
$$
\delta_k(x) = c_k^{-1}\left[\frac{\cos\left(\pi x\epsilon^{-1}\right)+1}{2}\right]^{(l+1)k}
$$
where $c_k = \int_{-\epsilon}^\epsilon \left[\frac{\cos\left(\pi s\epsilon^{-1}\right)+1}{2}\right]^{(l+1)k} \textnormal{d}s$. Note that $(\delta_k)$ is Dirac sequence and each $\delta_k$ is $l$ times continuous differentiable.


\begin{itemize}
\item[(i)] If there exists $t \in I$ such that $f(t) \not\in \RFA^0$, then $r(t) + q(t)a_m \neq 0$.  
Since $I$ is open, there exists $\epsilon > 0$ such that $(t-\epsilon, t+\epsilon) \subset I$ and $f$ is continuous in $(t-\epsilon, t+\epsilon)$. 
Using the functions $\delta_k$, we define the functions $\eta_k(z) = r_k(z) + q_k(z)A$, $\forall z \in [a,b]$, with 
$$
r_k(z) = \left\lbrace\begin{array}{cl}
r(z)\delta_k(z-t)    &, \mbox{ if } |z-t|< \epsilon \\
0     &, \mbox{ otherwise}  
\end{array}\right.
$$
and 
$$
q_k(z) = \left\lbrace\begin{array}{cl}
q(z)\delta_k(z-t)    &, \mbox{ if } |z-t|< \epsilon \\
0     &, \mbox{ otherwise}  
\end{array}\right.
$$
for all $z \in [a,b].$

Let $B_k = \int_{a}^{b} f(z)\odot\eta_k(z)\textnormal{d}z$ and $[B_k]_1 = \{b_k\}$ with
\begin{eqnarray*}
b_k & = & \left(\int_{a}^{b} r(z)r_k(z) - a_m^2q(z)q_k(z) \textnormal{d}z\right) +  \\ &  & \, +   
\left(\int_{a}^{b} r(z)q_k(z) + r_k(z)q(z) + 2a_mq(z)q_k(z) \textnormal{d}z\right)a_m \\
& = & \left(\int_{t-\epsilon}^{t+\epsilon} (r(z)^2 - a_m^2q(z)^2)\delta_k(z-t) \textnormal{d}z\right) +   \\ &  & \, +    
\left(\int_{t-\epsilon}^{t+\epsilon} (2r(z)q(z) + 2a_mq(z)^2)\delta_k(z-t) \textnormal{d}z\right)a_m \\
& = & \int_{t-\epsilon}^{t+\epsilon} (r(z)^2 + a_mq(z))^2\delta_k(z-t) \textnormal{d}z.
\end{eqnarray*}
Since $(\delta_k)$ is a Dirac sequence, we have that $b_k \to (r(t)^2 + a_mq(t))^2 > 0$ when $k \to \infty$. Thus, for $k$ sufficient large we have that $b_k > 0$, which implies that $B_k \not\in \RFA^0$. This contradicts the hypothesis that $\int_{a}^{b} f(t)\odot\eta_k(t)\textnormal{d}t \in \RFA^0$. 
Therefore, we must have that $f(t) \in \RFA^0$ almost everywhere. 

\item[(ii)] 
%Suppose there exists $t \in I$ such that $f(t) \neq 0$.  
%This implies that $r(t) \neq 0$ or $q(t) \neq 0$
For every $t \in I$, since $I$ is open, there exists $\epsilon > 0$ such that $(t-\epsilon, t+\epsilon) \subset I$ and $f$ is continuous in $(t-\epsilon, t+\epsilon)$. 
Using the functions $\delta_k$, we define the functions $\tilde{\eta}_k(z) = \tilde{r}_k(z)$, $\forall z \in [a,b]$, with 
$$
\tilde{r}_k(z) = \left\lbrace\begin{array}{cl}
\delta_k(z-t)    &, \mbox{ if } |z-t|< \epsilon \\
0     &, \mbox{ otherwise}  
\end{array}\right.$$
for all $z \in [a,b].$
From our hypothesis, we have
\begin{eqnarray*}
0 & = & \left(\int_{a}^{b} r(z)\tilde{r}_k(z) \textnormal{d}z\right) +     
\left(\int_{a}^{b} \tilde{r}_k(z)q(z)\textnormal{d}z\right)A \\
& = & \left(\int_{t-\epsilon}^{t+\epsilon} r(z)\delta_k(z-t) \textnormal{d}z\right) +     
\left(\int_{t-\epsilon}^{t+\epsilon} q(z)\delta_k(z-t)\textnormal{d}z\right)A.
\end{eqnarray*}
Using the above equality and the fact that $(\delta_k)$ is Dirac sequence, we obtain 
$$\Rightarrow \left\lbrace \begin{array}{l}
0 = \int_{t-\epsilon}^{t+\epsilon} r(z)\delta_k(z-t) \textnormal{d}z \xrightarrow[k\to \infty]{} r(t)     \\ 
0 = \int_{t-\epsilon}^{t+\epsilon} q(z)\delta_k(z-t) \textnormal{d}z \xrightarrow[k\to \infty]{} q(t)
\end{array}\right..
$$
%For $k$ sufficient large, this produces a contradiction with the hypothesis that $\int_{a}^{b} \tilde{\eta}_k(z)\odot f(z) \textnormal{d}z = 0$, since $r(t) \neq 0$ or $q(t) \neq 0$. 
Therefore, we conclude that $f(t) = 0$ almost everywhere in $[a,b].$
\end{itemize}
\end{proof}

\end{document}
